% ==================================================================
% Claude Code Workflow — AI-Driven Development in Practice
% Claude Code ワークフロー — AI駆動開発の実践
%
% This is a standalone slide deck, independent of the workshop slides.
% Delete this entire directory (docs/bonus/) to remove it.
% ==================================================================
\documentclass[aspectratio=169, 14pt]{beamer}

% --- Theme ---
\usetheme{Madrid}
\usecolortheme{default}

% --- Brand colors ---
\definecolor{sfblue}{RGB}{0, 120, 200}
\definecolor{sfdark}{RGB}{0, 60, 100}
\definecolor{sflight}{RGB}{220, 240, 255}
\definecolor{sfgray}{RGB}{80, 80, 80}
\definecolor{sfred}{RGB}{220, 50, 30}
\definecolor{sfgreen}{RGB}{40, 160, 40}

\setbeamercolor{palette primary}{bg=sfblue, fg=white}
\setbeamercolor{palette secondary}{bg=sfdark, fg=white}
\setbeamercolor{palette tertiary}{bg=sfdark, fg=white}
\setbeamercolor{structure}{fg=sfblue}
\setbeamercolor{title}{fg=white, bg=sfblue}
\setbeamercolor{frametitle}{fg=white, bg=sfblue}
\setbeamercolor{block title}{bg=sfblue!90, fg=white}
\setbeamercolor{block body}{bg=sflight}
\setbeamercolor{block title alerted}{bg=sfred!90, fg=white}
\setbeamercolor{block body alerted}{bg=sfred!8}
\setbeamercolor{block title example}{bg=sfgreen!80, fg=white}
\setbeamercolor{block body example}{bg=sfgreen!8}
\setbeamercolor{item}{fg=sfblue}

\setbeamerfont{title}{size=\Large, series=\bfseries}
\setbeamerfont{frametitle}{size=\large, series=\bfseries}

% --- Navigation & footer ---
\setbeamertemplate{navigation symbols}{}
\setbeamertemplate{footline}{%
  \leavevmode\hbox{%
    \begin{beamercolorbox}[wd=.33\paperwidth,ht=2.5ex,dp=1ex,center]{palette primary}%
      \usebeamerfont{author in head/foot}\insertshortauthor
    \end{beamercolorbox}%
    \begin{beamercolorbox}[wd=.34\paperwidth,ht=2.5ex,dp=1ex,center]{palette secondary}%
      \usebeamerfont{title in head/foot}\insertshorttitle
    \end{beamercolorbox}%
    \begin{beamercolorbox}[wd=.33\paperwidth,ht=2.5ex,dp=1ex,right]{palette tertiary}%
      \usebeamerfont{date in head/foot}%
      \insertframenumber{} / \inserttotalframenumber\hspace*{2ex}
    \end{beamercolorbox}%
  }\vskip0pt%
}

% --- Japanese support (LuaLaTeX) ---
\usepackage{luatexja}
\usepackage[match]{luatexja-fontspec}
\usepackage{ifthen}
\newboolean{jfontset}\setboolean{jfontset}{false}
\IfFontExistsTF{Hiragino Kaku Gothic ProN}{%
  \setmainjfont{Hiragino Kaku Gothic ProN}%
  \setsansjfont{Hiragino Kaku Gothic ProN}%
  \setboolean{jfontset}{true}}{}
\ifthenelse{\boolean{jfontset}}{}{%
  \IfFontExistsTF{Yu Gothic}{%
    \setmainjfont{Yu Gothic}%
    \setsansjfont{Yu Gothic}%
    \setboolean{jfontset}{true}}{}}
\ifthenelse{\boolean{jfontset}}{}{%
  \IfFontExistsTF{Noto Sans CJK JP}{%
    \setmainjfont{Noto Sans CJK JP}%
    \setsansjfont{Noto Sans CJK JP}}{}}

% --- Packages ---
\usepackage{listings}
\usepackage{url}
\usepackage{amsmath}
\usepackage{tikz}
\usetikzlibrary{arrows.meta, positioning, shapes.geometric, calc}

% --- Code listing style ---
\lstdefinestyle{cccode}{
  basicstyle=\ttfamily\tiny,
  keywordstyle=\color{sfblue}\bfseries,
  commentstyle=\color{sfgray}\itshape,
  stringstyle=\color{sfgreen},
  frame=single,
  framerule=0.5pt,
  rulecolor=\color{sfblue!40},
  backgroundcolor=\color{sflight},
  breaklines=true,
  showstringspaces=false,
  tabsize=2,
}
\lstset{style=cccode}

% --- Useful macros ---
\newcommand{\api}[1]{\texttt{\color{sfblue}#1}}
\newcommand{\warn}[1]{\textcolor{sfred}{\textbf{#1}}}

% === Document ===
\title{Claude Code ワークフロー}
\subtitle{AI駆動開発の実践}
\author{伊藤 恒平}
\date{2026}

\begin{document}

% === Title ===
\begin{frame}
  \titlepage
\end{frame}

% ------------------------------------------------------------------
\begin{frame}{このスライドについて / Purpose}
  \begin{block}{概要}
    StampFly Ecosystem のスライド・ファームウェア・ツールは
    すべて \textbf{Claude Code}（Anthropic の CLI エージェント）と人間の協働で開発された
  \end{block}

  \vspace{0.3em}

  \begin{itemize}
    \item 817 コミット / 70 日間 — アクティブ日平均 20 コミット
    \item fix 329 件、feat 225 件、docs 140 件
    \item CLAUDE.md の継続的改善で品質を担保
  \end{itemize}

  \vspace{0.3em}

  \begin{exampleblock}{ゴール}
    {\small 実プロジェクトで得た Claude Code のワークフローと Tips を共有する}
  \end{exampleblock}
\end{frame}

% ------------------------------------------------------------------
\begin{frame}{Claude Code とは / What is Claude Code?}
  \begin{block}{概要}
    Anthropic 公式の CLI エージェント。ターミナルから Claude を呼び出し、
    コード読み書き・ビルド・テスト・Git 操作を自律的に実行する
  \end{block}

  \vspace{0.3em}

  \begin{table}
    \centering\small
    \begin{tabular}{ll}
      \hline
      \textbf{機能} & \textbf{説明} \\
      \hline
      CLAUDE.md       & プロジェクト固有の指示書（常に読み込まれる） \\
      Memory          & セッション間で記憶を永続化（4タイプ） \\
      Plan Mode       & Shift+Tab で読み取り専用の設計モード \\
      Subagent        & 並列タスク・重い処理を子プロセスに委任 \\
      Hooks           & ツール実行前後にシェルコマンドを自動実行 \\
      Custom Skills   & \api{/commit} 等のカスタムワークフロー \\
      MCP Server      & 外部サービス連携（DB、draw.io 等） \\
      \hline
    \end{tabular}
  \end{table}

  {\footnotesize 公式: \url{https://docs.anthropic.com/en/docs/claude-code}}
\end{frame}

% ------------------------------------------------------------------
\begin{frame}{CLAUDE.md — プロジェクトの指示書}
  \begin{block}{階層構造（すべて同時に読み込まれる）}
    {\small
    \begin{tabular}{lp{60mm}}
      \api{\textasciitilde/.claude/CLAUDE.md} & グローバル設定（全プロジェクト共通）\\
      \api{<repo>/CLAUDE.md}      & プロジェクト設定（Git 管理） \\
      \api{<subdir>/CLAUDE.md}    & ディレクトリ固有の設定 \\
      \api{.claude/rules/*.md}    & パス限定ルール（paths フロントマター）\\
    \end{tabular}
    }
  \end{block}

  \vspace{0.2em}

  \begin{alertblock}{書くべき内容}
    \begin{itemize}\footnotesize
      \item ビルド・テストコマンド / コーディング規約 / コミット形式
      \item \warn{やってはいけないこと}（アンチパターン）← 最重要
      \item プロジェクト構造の概要
    \end{itemize}
  \end{alertblock}

  {\footnotesize \textbf{Tips:} 「ガードレール」型（禁止事項の明示）は指示型より効果的。200行以内を目標に}
\end{frame}

% ------------------------------------------------------------------
\begin{frame}{CLAUDE.md の実例 / Real Example}
  \begin{lstlisting}[language={}, numbers=none]
## Build Environment
sf build vehicle     # ビルド
sf flash vehicle -m  # 書き込み+モニタ

## Slide Rules
**CRITICAL: .tex 変更後は必ず PDF リビルド+レビュー**

## Slide Review Checklist
| ID | チェック内容                    |
|----|--------------------------------|
| L1 | フッターにはみ出していないか   |
| T4 | 矢印は水平・垂直のみ          |
| T8 | resizebox{w}{h} 両方指定は禁止 |
  \end{lstlisting}

  {\footnotesize チェックリストは実際の失敗事例から段階的に構築された（後述）}
\end{frame}

% ------------------------------------------------------------------
\begin{frame}{Memory — セッション間の記憶}
  \begin{columns}[T]
    \begin{column}{0.48\textwidth}
      \begin{block}{4つのメモリタイプ}
        {\small
        \begin{tabular}{lp{28mm}}
          \hline
          \textbf{Type} & \textbf{用途} \\
          \hline
          user     & ユーザーの専門性・好み \\
          feedback & 修正指示・NG パターン \\
          project  & 進行中の作業・期限 \\
          reference & 外部リソースへのリンク \\
          \hline
        \end{tabular}
        }
      \end{block}
    \end{column}
    \begin{column}{0.48\textwidth}
      \begin{exampleblock}{feedback メモリの実例}
        {\scriptsize
        \texttt{速度優先は手戻りを招き工数増大。}\\
        \texttt{チェックリストは1回もスキップしない。}\\
        \texttt{サブエージェントの判定を鵜呑みにしない。}
        }
      \end{exampleblock}
      \vspace{0.3em}
      {\footnotesize 一度保存すれば、\textbf{全ての将来のセッション}で Claude がこのルールに従う}
    \end{column}
  \end{columns}

  \vspace{0.5em}

  {\footnotesize \textbf{Tips:} コードや Git 履歴から分かる情報は保存しない（二重管理の回避）}
\end{frame}

% ------------------------------------------------------------------
\begin{frame}{Plan Mode — 実装前に設計する}
  \begin{block}{Plan Mode とは}
    Shift+Tab でトグル。コード編集を\textbf{できない}モード。\\
    調査・設計に集中させ、方針が固まってから実装に移る
  \end{block}

  \vspace{0.2em}

  \begin{columns}[T]
    \begin{column}{0.45\textwidth}
      \begin{alertblock}{有効な場面}
        \begin{itemize}\small
          \item 複数ファイルにまたがる変更
          \item アーキテクチャの判断が必要
          \item 失敗コストが高い作業
        \end{itemize}
      \end{alertblock}
    \end{column}
    \begin{column}{0.5\textwidth}
      \begin{exampleblock}{典型的なフロー}
        {\small
        \textbf{1.} Plan Mode で調査・設計 \\
        \textbf{2.} ユーザーが計画を承認 \\
        \textbf{3.} 実装 → テスト → コミット \\
        \textbf{4.} 問題発見 → 計画を更新
        }
      \end{exampleblock}
    \end{column}
  \end{columns}

  \vspace{0.3em}

  {\footnotesize \textbf{Tips:} 「思考」と「実行」を分離することで手戻りを大幅に削減できる}
\end{frame}

% ------------------------------------------------------------------
\begin{frame}{Subagent — 並列処理とコンテキスト保護}
  \begin{table}
    \centering\footnotesize
    \begin{tabular}{llp{48mm}}
      \hline
      \textbf{タイプ} & \textbf{用途} & \textbf{説明} \\
      \hline
      general-purpose & 汎用タスク     & 全ツール利用可。複雑な作業を委任 \\
      Explore         & コード検索     & 読み取り専用。高速なコードベース探索 \\
      Plan            & 設計           & 読み取り専用。アーキテクチャ設計 \\
      \hline
    \end{tabular}
  \end{table}

  \begin{exampleblock}{このプロジェクトでの活用例}
    \begin{itemize}\footnotesize
      \item \textbf{画像レビュー} — PDF を画像化して目視確認。メインに画像を持ち込まない
      \item \textbf{並列リサーチ} — 複数サブエージェントを同時起動して情報収集
      \item \textbf{バックグラウンド実行} — ビルドを裏で実行し完了通知を受け取る
    \end{itemize}
  \end{exampleblock}

  \begin{alertblock}{}
    {\footnotesize サブエージェントの「問題なし」を鵜呑みにしない — メイン側で検証すること}
  \end{alertblock}
\end{frame}

% ------------------------------------------------------------------
\begin{frame}{Hooks \& Skills — 自動化の仕組み}
  \begin{columns}[T]
    \begin{column}{0.48\textwidth}
      \begin{block}{Hooks（自動実行）}
        {\small
        \begin{tabular}{lp{26mm}}
          \hline
          \textbf{Event} & \textbf{用途} \\
          \hline
          PreToolUse  & 危険操作のブロック \\
          PostToolUse & 自動 lint / format \\
          Notification & Slack 通知転送 \\
          Stop        & 応答完了後の自動テスト \\
          \hline
        \end{tabular}
        }

        \vspace{0.3em}

        {\scriptsize \api{settings.json} で設定。\\
        prompt-based / agent-based も可能}
      \end{block}
    \end{column}
    \begin{column}{0.48\textwidth}
      \begin{block}{Custom Skills}
        {\small
        \api{/commit} — ガイドライン準拠の\\
        \hspace{2em}Git コミットを自動作成 \\[4pt]
        \api{/simplify} — コード品質チェック \\[4pt]
        \api{/batch} — 並列ワークツリーで\\
        \hspace{2em}大規模変更を一括処理
        }

        \vspace{0.3em}

        {\scriptsize \api{.claude/skills/} に SKILL.md で定義。\\
        \texttt{\$ARGUMENTS} で引数を受け取れる}
      \end{block}
    \end{column}
  \end{columns}

  \vspace{0.3em}

  {\footnotesize \textbf{Tips:} 繰り返す作業は Skill に、条件付き自動化は Hook にまとめる}
\end{frame}

% ------------------------------------------------------------------
\begin{frame}{コンテキスト管理 / Context Management}
  \begin{alertblock}{コンテキストウィンドウは最重要リソース}
    Claude が一度に扱える情報量には上限。無駄に消費すると性能劣化し指示を忘れる
  \end{alertblock}

  \vspace{0.2em}

  \begin{table}
    \centering\small
    \begin{tabular}{ll}
      \hline
      \textbf{手法} & \textbf{説明} \\
      \hline
      \api{/compact}       & 会話履歴を要約して圧縮（\textasciitilde95\% で自動発動）\\
      Subagent 委任         & 大量の検索結果をメインに持ち込まない \\
      Memory 永続化         & 繰り返し伝える情報はメモリに保存 \\
      CLAUDE.md 簡潔化      & 毎回読み込まれるので 200 行以内に \\
      画像はサブエージェント & PDF/画像の Read はコンテキスト消費大 \\
      \hline
    \end{tabular}
  \end{table}

  {\footnotesize \textbf{教訓:} PDF 全ページ確認をメインで行い rate limit 抵触 → サブエージェント限定ルールを制定}
\end{frame}

% ------------------------------------------------------------------
\begin{frame}{実践ワークフロー / Workflow in Practice}
  \begin{center}
    \begin{tikzpicture}[
      node distance=6mm and 10mm,
      box/.style={rectangle, rounded corners=3pt, draw=sfblue,
                  fill=sflight, minimum height=8mm, minimum width=28mm,
                  font=\small, align=center},
      arr/.style={-{Stealth[length=2.5mm]}, sfblue, thick},
    ]
      \node[box] (start)   {前回の Next steps\\を確認};
      \node[box, right=of start] (plan) {Plan Mode\\で設計};
      \node[box, right=of plan]  (impl) {実装\\（小さな単位）};
      \node[box, below=of impl]  (test) {ビルド/テスト\\で検証};
      \node[box, below=of plan]  (commit) {\texttt{/commit}\\で記録};
      \node[box, below=of start] (review) {レビュー\\（チェックリスト）};

      \draw[arr] (start)  -- (plan);
      \draw[arr] (plan)   -- (impl);
      \draw[arr] (impl)   -- (test);
      \draw[arr] (test)   -- (commit);
      \draw[arr] (commit) -- (review);

      % Fix loop
      \draw[arr, sfred] (review.west) -- ++(-6mm,0)
        node[above, font=\scriptsize, sfred, pos=0.4] {問題あり}
        |- (impl.west);

      % Done
      \draw[arr, sfgreen] (review.south) -- ++(0,-6mm)
        node[right, font=\scriptsize, sfgreen] {完了 → 次のタスクへ};
    \end{tikzpicture}
  \end{center}
\end{frame}

% ------------------------------------------------------------------
\begin{frame}{Git コミット戦略 / Commit Strategy}
  \begin{block}{Conventional Commits + Next steps}
    \begin{lstlisting}[language={}, numbers=none]
fix(slides): fix P33 arrow-node penetration

Changes:
- TikZ座標を調整しノード貫通を解消
- adjustbox で均等スケーリングに変更

Next steps:
- sf build vehicle でビルド確認
- P34-P40 の同様の問題を確認
    \end{lstlisting}
  \end{block}

  \vspace{0.3em}

  \begin{columns}[T]
    \begin{column}{0.45\textwidth}
      {\small \textbf{自動コミットのトリガー:}}
      \begin{itemize}\small
        \item ファイル変更後
        \item テスト/ビルド成功後
        \item 「ここまで」「休憩」等の発言
      \end{itemize}
    \end{column}
    \begin{column}{0.5\textwidth}
      {\small \textbf{Next steps の効果:}}
      \begin{itemize}\small
        \item セッション再開時の起点になる
        \item 817 コミット中 318 件に記載
        \item 具体的なコマンド・検証項目を含む
      \end{itemize}
    \end{column}
  \end{columns}
\end{frame}

% ------------------------------------------------------------------
\begin{frame}{失敗から学ぶ / Lessons from Failures}
  \begin{alertblock}{事例: P33 espnow\_dataflow — 7回の fix コミット}
    1枚のスライドの TikZ 図修正にリビルドサイクルが12回以上発生
  \end{alertblock}

  \vspace{0.3em}

  \begin{columns}[T]
    \begin{column}{0.48\textwidth}
      {\small \textbf{何が起きたか:}}
      \begin{enumerate}\scriptsize
        \item \texttt{resizebox\{w\}\{h\}} でフォント縦伸び
        \item 矢印方向が上段/下段で不統一
        \item 場当たり的に Beamer 側で対処
        \item チェックリスト未適用で見落とし
      \end{enumerate}
    \end{column}
    \begin{column}{0.48\textwidth}
      {\small \textbf{得られた教訓:}}
      \begin{enumerate}\scriptsize
        \item 根本原因を最初に特定する
        \item チェックリストは1回も省略しない
        \item サブエージェントの判定を鵜呑みにしない
        \item CLAUDE.md に失敗事例を記録 → 再発防止
      \end{enumerate}
    \end{column}
  \end{columns}

  \vspace{0.5em}

  \begin{exampleblock}{対策}
    {\small feedback メモリに保存 → 全セッションで Claude がこのルールに従うようになった}
  \end{exampleblock}
\end{frame}

% ------------------------------------------------------------------
\begin{frame}{反復改善のパターン / Iteration Patterns}
  \begin{table}
    \centering\small
    \begin{tabular}{lcp{48mm}}
      \hline
      \textbf{領域} & \textbf{Fix} & \textbf{内容} \\
      \hline
      installer          & 18 & WSL2、Windows パス、Python 検出 \\
      シミュレータカメラ   & 13 & VPython カメラ追従（revert 2回含む） \\
      P33 TikZ 図        &  7 & ルーティング、アスペクト比、ノード貫通 \\
      L10/L11 ヘッダー   &  5 & framesubtitle 削除後のレイアウト崩れ \\
      エコシステム図      &  5 & 矢印ラベル重なり → 直角ルーティング \\
      コントローラメニュー &  4 & フリッカー、残像、スクロール \\
      \hline
    \end{tabular}
  \end{table}

  \vspace{0.5em}

  \begin{block}{パターンの分析}
    \begin{itemize}
      \item \textbf{視覚的な問題}は反復回数が多い（テキストレビューだけでは不十分）
      \item \textbf{チェックリスト導入後}は反復回数が減少
      \item revert は 817 コミット中わずか 10 件（1.2\%）→ 方向転換は稀
    \end{itemize}
  \end{block}
\end{frame}

% ------------------------------------------------------------------
\begin{frame}{CLAUDE.md の進化 / Evolution of CLAUDE.md}
  \begin{center}
    \begin{tikzpicture}[
      node distance=3.5mm,
      evt/.style={rectangle, rounded corners=2pt, draw=sfblue,
                  fill=sflight, minimum width=80mm, minimum height=5.5mm,
                  font=\scriptsize, align=left},
      arr/.style={-{Stealth[length=2mm]}, sfgray},
    ]
      \node[evt] (e1) {Jan: ビルドコマンド、コーディング規約、sf CLI 使用方針};
      \node[evt, below=of e1] (e2) {Jan: コミットガイドライン + Next steps ルール追加};
      \node[evt, below=of e2] (e3) {Feb: スライドレビューチェックリスト導入（L1--L6, T1--T9）};
      \node[evt, below=of e3] (e4) {Feb: サブエージェントによる画像レビュー方式に切替};
      \node[evt, below=of e4] (e5) {Mar: サブエージェント限定ルール強制化（rate limit 対策）};
      \node[evt, below=of e5] (e6) {Mar: TikZ チェックリスト大幅拡充（T4 直角ルーティング等）};
      \node[evt, below=of e6] (e7) {Mar: feedback メモリで品質優先の原則を永続化};

      \draw[arr] (e1) -- (e2);
      \draw[arr] (e2) -- (e3);
      \draw[arr] (e3) -- (e4);
      \draw[arr] (e4) -- (e5);
      \draw[arr] (e5) -- (e6);
      \draw[arr] (e6) -- (e7);
    \end{tikzpicture}
  \end{center}
\end{frame}

% ------------------------------------------------------------------
\begin{frame}{統計で見るプロジェクト / Project by Numbers}
  \begin{table}
    \centering
    \begin{tabular}{lr}
      \hline
      \textbf{指標} & \textbf{値} \\
      \hline
      総コミット数              & 817 \\
      開発期間                  & 70 日間（2026/1/5 -- 3/15）\\
      アクティブ日の平均         & 20.4 コミット/日 \\
      最多の日                  & 63 コミット（1/12）\\
      fix 比率                  & 40.3\%（329/817）\\
      feat 比率                 & 27.5\%（225/817）\\
      最長 fix チェイン          & 18 回（installer）\\
      revert 回数               & 10 回（1.2\%）\\
      Next steps 記載率         & 38.9\%（318/817）\\
      CLAUDE.md 改訂回数        & 17 回 \\
      \hline
    \end{tabular}
  \end{table}
\end{frame}

% ------------------------------------------------------------------
\begin{frame}{Top 10 Tips}
  \begin{enumerate}\small
    \item \textbf{CLAUDE.md は基盤} — ビルドコマンド、規約、禁止事項を書く
    \item \textbf{ガードレール $>$ 指示} — 「やるな」の方が「やれ」より効果的
    \item \textbf{Plan → Implement → Test} — 考えてから書く
    \item \textbf{小さくコミット} — 1機能 = 1コミット、Next steps を含める
    \item \textbf{チェックリストは省略しない} — 速度優先は手戻りで工数増大
    \item \textbf{feedback メモリで学習} — 同じミスを二度と繰り返さない
    \item \textbf{サブエージェントで並列化} — コンテキストを守りつつ効率化
    \item \textbf{画像レビューはサブエージェントに} — rate limit 対策
    \item \textbf{CLAUDE.md を育てる} — 失敗するたびにルールを追加
    \item \textbf{コンテキストは最重要リソース} — \api{/compact} と Memory を活用
  \end{enumerate}
\end{frame}

% ------------------------------------------------------------------
\begin{frame}{参考資料 / References}
  \begin{table}
    \centering\small
    \begin{tabular}{lp{55mm}}
      \hline
      \textbf{リソース} & \textbf{URL} \\
      \hline
      公式ドキュメント &
        \url{https://docs.anthropic.com/en/docs/claude-code} \\
      ベストプラクティス &
        \url{https://docs.anthropic.com/en/docs/claude-code/best-practices} \\
      Claude Code GitHub &
        \url{https://github.com/anthropics/claude-code} \\
      \hline
    \end{tabular}
  \end{table}

  \vspace{0.5em}

  \begin{exampleblock}{Claude Code を始めるには}
    {\small
    \texttt{npm install -g @anthropic-ai/claude-code} \\
    \texttt{cd your-project \&\& claude}
    }
  \end{exampleblock}

  \vspace{0.3em}

  \begin{block}{インストール後の最初の一歩}
    {\small
    \texttt{claude /init} — CLAUDE.md を自動生成してからカスタマイズ
    }
  \end{block}
\end{frame}

\end{document}
